% Generated draft: requires mathematical and source review.
\documentclass[11pt]{article}
\usepackage[margin=25mm]{geometry}
\usepackage{fontspec}
\setmainfont{Latin Modern Roman}
\usepackage{amsmath,amssymb,mathtools,unicode-math}
\setmathfont{Latin Modern Math}
\usepackage{enumitem,needspace,xurl,hyperref}
\setlength{\parindent}{0pt}
\setlength{\parskip}{5pt}
\setlength{\emergencystretch}{3em}
\newcommand{\N}{\mathbb{N}}
\newcommand{\Z}{\mathbb{Z}}
\newcommand{\Q}{\mathbb{Q}}
\newcommand{\R}{\mathbb{R}}
\newcommand{\C}{\mathbb{C}}
\newcommand{\F}{\mathbb{F}}
\newcommand{\A}{\mathbb{A}}
\newcommand{\PP}{\mathbb{P}}
\newcommand{\E}{\mathbb{E}}
\newcommand{\eps}{\varepsilon}
\newcommand{\dd}{\,\mathrm{d}}
\begin{document}
\section*{UnsolvedMath ID: 30002036}
\textbf{Source problem code:} catalogue entry
\subsection{Definitions and mathematical statement}
To define the notation, let \(k\) be a field, let \(n\geq 0\) and \(d\geq 1\) be integers, and put \(V=H^0(\PP^n_k,\mathcal O_{\PP^n_k}(d))\). The complete linear system defines the \(d\)-uple Veronese embedding
\[
v_d\colon \PP^n_k\longrightarrow \PP(V^*)\cong\PP^{\binom{n+d}{d}-1}_k.
\]
Let
\[
S=\operatorname{Sym}_k(V),\qquad
R=\bigoplus_{m\geq 0}H^0(\PP^n_k,\mathcal O_{\PP^n_k}(md)),
\]
where \(V\) has degree one in \(S\), and the summand indexed by \(m\) has degree \(m\) in \(R\). Multiplication of sections gives the natural graded map \(S\to R\). For this embedding, property \(N_p\), with \(p\geq 0\), means that this map is surjective and that the minimal graded free resolution of \(R\) over \(S\) has the form
\[
\cdots\longrightarrow F_2\longrightarrow F_1\longrightarrow S\longrightarrow R\longrightarrow 0,
\qquad F_i\cong S(-i-1)^{\beta_i}\quad(1\leq i\leq p),
\]
for nonnegative integers \(\beta_i\). Here \(S(-a)\) denotes the graded free module whose generator has degree \(a\); zero modules are allowed. Thus \(N_0\) requires projective normality, while positive \(p\) additionally requires linearity of the indicated syzygies.

The supplied corrected question asks whether, throughout its intended parameter range,
\[
v_d(\PP^n_k)\text{ satisfies property }N_{3d-3}.
\]
The source record does not specify the ground field, its characteristic, or precise bounds on \(n\) and \(d\). The ranges used above make the notation meaningful; they do not establish that every such choice is intended by the source. In particular, no characteristic-free universal assertion is attributed to the report here.

The original extracted statement also asked how to determine efficiently an implicit equation from approximation complexes and multigraded syzygies, given a divisor \(D\) on \(X\) and a basepoint-free subspace \(U\subset H^0(X,\mathcal O_X(D))\) with \(\dim U=\dim X+2\) and hypersurface image in the source's notation \(\PP(U)\). The supplied correction explicitly assigns that material to the next record, so it is retained here only as extraction history, not as a second mathematical target. Its hypotheses on \(X\), projectivization convention, and meaning of efficiency are not specified.
\subsection{Short English statement}
The corrected question asks whether the degree-\(d\) Veronese embedding has property \(N_{3d-3}\): projective normality together with linear syzygies through homological degree \(3d-3\). The intended field, characteristic, and parameter range require clarification. The implicitization question in the original extraction is assigned to the next record by the supplied correction.
\subsection{Sources}
\begin{itemize}
\item Local UnsolvedMath registry: permanent ID \(30002036\), source problem code \texttt{catalogue entry}, title ``Property \(N_{3d-3}\) for Veronese Embeddings,'' in the supplied Oberwolfach Reports Open Problems collection. No individual proposer is supplied; the proposed year is recorded as 2012.
\item Supplied report citation: \emph{Toric Geometry (2013)}, Oberwolfach Reports, pp.\ 1231--1306, DOI \texttt{10.4171/owr/2012/21}: \url{https://doi.org/10.4171/owr/2012/21}. The title year 2013 is retained alongside the DOI year component 2012 and the recorded proposed year 2012; their relationship has not been verified.
\item Supplied statement-verification note: the official report poses the Veronese-syzygy question separately, and the extracted implicitization material belongs to the next record. The registry labels the statement \texttt{corrected\_verified}; this is supplied curation, not an independent verification performed for this draft.
\item Supplied literature reference, labelled ``Ottaviani--Paoletti'': \url{https://arxiv.org/abs/math/9811131}.
\item Supplied literature reference, labelled ``Vu on third Veronese embeddings'': \url{https://arxiv.org/abs/1303.5532}.
\item Supplied literature reference, labelled ``Recent syzygy implicitization'': \url{https://arxiv.org/abs/2606.01415}. This reference accompanies the implicitization material separated from the corrected target.
\item Literature assessment checked at \texttt{2026-08-21}: the supplied registry labels the record partially solved and states that the Veronese syzygy conjecture is proved for degree three and several special cases but remains open in general. It also reports substantial approximation-complex and syzygy progress for toric hypersurface implicitization without a complete general theory. This is supplied dated curation derived from the OWR collection, not independently verified literature assessment; the source itself requests independent verification of time-sensitive status and cited claims.
\end{itemize}
\end{document}
